% Materials and Methods, wrapped around the five figures (rebuilt 2026-09-24 against the figure builds of that
% morning; the sixth figure, the LevelCheck comparison, moved to paper 2 that evening). Every number is traced in
% NUMBERS.md ("The six-figure wrap" and "Revision 2") and sections/methods_sources.md.
% Order here = order of the Results, with the same subheadings (REVIEW2_RSNA.md Fix 25): population and sources,
% phenotypes (Figure 1, Table 2), the reference standard (the automated reference of LEVELCHECK_MODEL.md s14 moved
% to paper 2 in the 2026-09-24 fill pass: no paper-1 result uses it; paper2_levelcheck/MOVED_FROM_PAPER1.md item 8),
% the encoder and its embedding (Figure 2), the network (Figure 3), naming, counting
% and coupling on the intraoperative C-arm (Figure 4) and transitional sample size (Figure 5). First paragraph:
% design, dates, ethics, subject overlap. Last paragraph: statistics. Device-card formation, the corpus manifold,
% the detector's source contributions and detection against angle are in the supplement (Figures S2-S5, Tables
% S1-S4).
% Fill pass (2026-09-24 evening, REVISION_LOG3.md section 8): "Design" dates = April 2026 (month precision: the
% TCIA CT COLONOGRAPHY download scripts are in the repository's first commit, 550cfff, 2026-04-18; no on-grid file of
% that download survives, so the day is not defended) to 24 September 2026. The two CTSpinoPelvic1K readers are
% radiology residents (Greg 2026-09-24 22:15: both PGY2; earlier that day names only, no year was recorded in either repository, so
% none is printed; the descriptor's title page carries their names and department, which the anonymized file may not).
% Revision 5 (2026-09-24, REVISION_LOG3.md section 8; PAPER1_SCOPE.md): the paragraph "The LevelCheck comparison"
% (comparator, strata, unit, s15 detector, s12-s13 secondary analyses), the power sentence of "Statistical
% analysis" and the one-sided .025 clause moved verbatim to paper2_levelcheck/MOVED_FROM_PAPER1.md item 5. What
% stays of that material is paper 1's own cross-validation reporting: the coupled call, the count and the stop
% signal ("The network"; Table 3, Results). "Transitional sample size" is rewritten around the pre-stated Wilson
% targets and defines "at risk" here, since Table 3 no longer does.
% Revision 2 (2026-09-24, manuscript-drive; REVISION_LOG2.md): EVAL_LEAKS.md applied (the zero-shot run is the
% lineage-complete v9 encoder and head, HELD OUT for all 148 cases; the fold models are being re-run from it, their
% numbers PENDING G5); REVIEW2_RSNA.md Fixes 6, 11, 14, 16, 19, 20, 21, 25, 26 and cuts CUT2, CUT5, CUT6, CUT8;
% REVIEW2_STATS.md N1, N4, N12, N13, N17; LEVELCHECK_MODEL.md s14 (automated reference; s15, the primary detector,
% left with the comparison). The automated-reference selection rule is committed (b511985), the rule itself is
% frozen with its SHA-256 before scoring (s14 R2).
% Revision 3 (2026-09-24, collect-and-submit-1; REVISION_LOG3.md): one sentence at the end of "Statistical analysis"
% points to the supplement section "Data integrity and quality controls" (the evaluation gate and every audit tool);
% the abstention rule moved to the naming table's note d for room; small trims measured with check_limits.words().
\section*{Materials and Methods}

\paragraph{Design.}
This retrospective study used de-identified image collections (Table~\ref{tab:data}), analysed
from April 2026, when the first collection was downloaded, to 24 September 2026; intraoperative exposures carry study dates
11 November--12 December 2023. The institutional human research protection office determined this secondary
analysis not to be human participant research (determination 2026-269); consent and the Health Insurance
Portability and Accountability Act do not apply. Every subject was reported by its release or article; CTSpinoPelvic1K also in a
data descriptor \selfcite{ctsp1kdescriptor}.

\paragraph{Sources and transitional phenotypes.}
Rendered CT came from CTSpinoPelvic1K \selfcite{ctsp1k} (801 of 802 records), TotalSegmentator v2 and v3
\citep{totalsegmentator} and VerSe 19/20 \citep{verse2021}; CSXA \citep{csxa}, BUU-LSPINE \citep{buu2023} and
NHANES~II \citep{nhanes} supplied radiographs with released levels or landmarks, and VinDr-SpineXR pseudo-labelled
films \citep{vindr}. Real exposures with three-dimensional truth came from C-ArmLSpine \citep{carmlspine} (296
exposures, 148 of 149 cases, one mobile C-arm, per-case cone-beam CT), DeepFluoro (six specimens
\citep{grupp2020}), COBRA2026 (145 of 467 Elekta cases; seven of 152 downloaded failed a pre-export check
\citep{cobra2026}), a phantom \citep{genevahip} and a cadaveric specimen \citep{ljubljana2d3d}
(Table~\ref{tab:data}; flow, Figure~S1). Transitional persons were typed under a rib
convention as L6, T13, sacralised L5
(Castellvi III or IV \citep{castellvi1984}), Castellvi I/II or lumbar rib (Figure~\ref{fig:phenotypes},
Table~\ref{tab:transitional}): VerSe from its published read; CTSpinoPelvic1K from two radiology residents
(postgraduate year 2) grading independently, six disagreements resolved by consensus
\selfcite{ctsp1kdescriptor}; TotalSegmentator by an independent rib count; intraoperative and COBRA2026 calls are an
automatic labeller's (SPINEPS with VERIDAH \citep{veridah2026}), read by no person.

\paragraph{Reference standard.}
No person read an intraoperative image. Truth is the per-case cone-beam CT, registered to each exposure and
labelled by TotalSegmentator and by SPINEPS with VERIDAH; a vertebra is scored only where both agree (814 of
910), withholding T11 and most of T12, and both count from the sacrum, so agreement is not independent. Naming matches detections one to one to vertebral-body truth boxes
(Table~\ref{tab:naming}); detection is scored against whole-vertebra projections.

\paragraph{The level encoder and what it learns.}
Each body is read from a fixed-scale two-channel crop (image and target marker; target 96 of 320 pixels) by a
ResNet-18 encoder \citep{he2016} into a 512-dimensional embedding. The latest encoder's embedding was mapped by
t-distributed stochastic neighbour embedding (t-SNE) on its validation split after dropping 2,889 crops of 164
persons also in its training keys (Figure~\ref{fig:embedding}); that split chose the checkpoint, so it is not
held out. Separation from the level above is the cosine silhouette, bootstrapped by person and repeated per
render source. The three contested C-ArmLSpine L6 patients, in neither split, and the last free vertebra
of validation persons with an unlabelled fused L5 were placed against the neighbouring clusters (supplement).
Occlusion over ten lateral crops per level from persons outside the training keys gave the share and density
of attribution on the body above, the target and the body below. Renders were styled by image-formation terms
fitted per device, which missed their pre-stated criterion (Figures~S2 and S3, Table~S1).

\paragraph{The network.}
Two detectors never trained on the test C-arm find the bodies (Figure~\ref{fig:network}): a YOLO11m-pose keypoint
detector trained on films boxes the bodies the level model names (confidence 0.35), and a YOLO11m vertebra and
sacrum detector, fine-tuned on real frames of two other devices, boxes the sacrum at confidence 0.74, chosen on
three other devices (Table~S2); 97\,\% detection was pre-stated as counting grade. A two-layer
bidirectional gated recurrent unit (GRU) decodes runs of two to nine embeddings at mixed lengths into a 27-level
posterior. Counting ranks each body from the detected sacrum and reads no morphology. The coupled call chooses,
among the ordinary, L6, T13 and both-variant naming alphabets, the one best fitting the posteriors at the
ordinals; a classifier name differing from the call, or an alphabet margin below the fifth percentile of correct
frames in the other folds, stops the read (Table~\ref{tab:naming}).

\paragraph{Naming, counting and coupling on the intraoperative C-arm.}
C-ArmLSpine was scored zero-shot by the latest encoder and its recurrent head, whose training lists hold no
C-ArmLSpine case; with five case-grouped folds; and, because a case is an
acquisition session rather than a proven person (five dates), with three folds by acquisition day. Per
fold that encoder and head were fine-tuned (15 epochs, learning rate $10^{-4}$), checkpoints chosen on validation
folds excluding every test case (Figure~\ref{fig:exposure}). The primary outcome, designated after a first
zero-shot run was scored, was the exact level of detected lateral vertebrae at full coverage; secondary outcomes
were naming with one or three bodies in view, the count, the coupled call and its stop signal, and detection on
the 56 test frames (Table~S2, Figure~S4). Same-device
detection used nested whole-case subsets of 0 to 210 real frames, and one view alone at 25 and 50 frames (two
seeds; supplement). Detection against view angle used $\pm$90$^\circ$ render pose pools per
device and one patient's 301-frame rotational acquisition, the real pool holding only anteroposterior and lateral
views (Figure~S5). Naming against training persons used NHANES~II, BUU-LSPINE and CSXA with a render-only
encoder on nested person subsets, and C-ArmLSpine per view by case-grouped folds. Render-only detector arms, the fit of accuracy against count, and the
ablation of the rendered rows are in the supplement.

\paragraph{Transitional sample size.}
Pre-stated Wilson targets, a 95\,\% lower bound of 90\,\% for the exact level and 85\,\% for variant sensitivity on
pooled at-risk test persons, set the persons a transitional test needs (Table~\ref{tab:transitional},
Figure~\ref{fig:power}); at risk are the subtypes a count from the sacrum or from C2 misnames: L6, T13 and a
sacralised L5 on a four-lumbar count. The person-disjoint split (50:10:40 per subtype) puts every real transitional
subject on the test side; of its transitional test persons only the three C-ArmLSpine L6 patients, scored in
cross-validation, enter a result here.

\paragraph{Statistical analysis.}
Naming used every available case. Transitional sample size is the smallest person count at which the Wilson
95\,\% lower bound with zero to five errors reaches a target (Figure~\ref{fig:power}a). Intervals are 95\,\% two-thousand-resample bootstraps clustered by case,
patient or person, percentile with the bias-corrected and accelerated (BCa) interval beside; where errors fall in
few clusters, the Wilson interval on the design-effect-adjusted count replaces it. Paired comparisons report the same-item difference with its bootstrap interval and an exact
sign test over clusters, Holm-adjusted over four comparisons named before adjustment (lateral versus
anteroposterior naming; three in the supplement); other contrasts are exploratory; significance was two-sided
$P<.05$. Every reported number passed a patient-keyed evaluation gate
that counts a model's, and its ancestors', training and checkpoint-selection persons as seen (supplement, with
the other integrity checks). Analyses used Python 3.11 and 3.12 (Python Software Foundation,
Wilmington, Del), PyTorch 2.6.0 (PyTorch Foundation), Ultralytics 8.4 (Ultralytics, Frederick, Md), NumPy 2.2.6
and SciPy 1.15.3 and 1.18.1; reporting follows CLAIM 2024 \citep{claim2024}.
